\section{Introduction}
\label{sec:introduction}

Digital payments in East Africa run largely over mobile money: person-to-person transfers,
agent cash-in and cash-out, merchant payments and USSD sessions on basic handsets, alongside
cards and bank transfers. The fraud that accompanies them (SIM-swap account takeover, agent float
manipulation, mule accounts, social-engineering scams) is shaped by those channels. Much of the
published practice of fraud detection was developed on card transactions
\cite{bolton2002fraud,dalpozzolo2014lessons}, while mobile money fraud takes forms such as
SIM-swap takeover \cite{lee2020simswap}, attacks through add-on SIM hardware
\cite{phipps2018thinsim} and scams carried out by SMS and phone call
\cite{razaq2021scams,dogan2025grift}, on channels such as USSD that carry no device signals. We do not
repeat the frequently quoted figures about the growth of mobile money fraud in the region: the
claims register of this project lists each of them as unverified or replaced, and this paper
states none of them as fact (Section~\ref{sec:limitations}).

FraudShield is a fraud detection system built to a written specification (a software
requirements specification, SRS, corrected by a register of 51 binding resolutions) for this
setting: a synchronous decision path that returns a decision to core banking with a server-side
latency target, an explainable scoring model, and an analyst console. Building it required data
on which the model could be trained and evaluated. No public, labelled, East African
transaction dataset was available to us, so we built one: FraudShield-EAC, an agent-based,
seeded simulator of customers, agents, merchants and eight fraud scenarios across Rwanda, Kenya,
Tanzania, Uganda and the Democratic Republic of the Congo.

A synthetic benchmark invites a particular failure. It is easy to generate one on which a model
scores highly and hard to know whether the score measures learning or the generator. Much of this
project's effort went into interrogating the benchmark until it gave up its structure, and
publishing each interrogation whether or not it flattered the dataset. That record, not the
benchmark, is the first contribution.

\paragraph{Contributions.} We claim four things, each of which a reader can check against the
repository and disagree with:
\begin{enumerate}
  \item \textbf{A benchmark-construction methodology and its measured failure modes}
        (Sections~\ref{sec:dataset}, \ref{sec:evaluation} and \ref{sec:discussion}). A single
        feature reaches \nFloorM{}, so a model AUC quoted against chance overstates the quantity
        that matters; four disjoint feature groups each nearly solve the task alone; and the
        record of controls that failed, how each failure was found, and what now enforces it.
  \item \textbf{A detection failure on a pre-registered non-burst variant}
        (Section~\ref{sec:generalisation}): recall \nRevRecall{} \nRevCI{} against
        \nBaseVarRecall{} \nBaseVarCI{}, measured on a held-out shape rather than inferred from
        ablation margins.
  \item \textbf{Two generalisation nulls, and why this benchmark could not make them
        informative} (Section~\ref{sec:generalisation}): leave-one-country-out and a novel
        variant held out in time both fail to test generalisation, because the benchmark
        encodes fraud as bursts and its fraud mechanisms are country-invariant by construction.
  \item \textbf{The latency--explainability frontier} (Section~\ref{sec:evaluation}): the cost
        of an exact TreeSHAP explanation across tree complexity, reported at the median only,
        because the development machine cannot produce a repeatable 99th percentile.
\end{enumerate}
The redundancy finding (item 1) is stated separately in Section~\ref{sec:evaluation} because it
is a result about how scenario-driven generators encode signal, not a detail of one table.

The system itself, \nFeatures{} features computed on two independent paths with a parity suite, a
published temporal split, and a governance apparatus that refuses stale evidence, is the setting
of these results and is not claimed as a research contribution. Its deviations from the
specification are recorded as architecture decision records (ADRs), which we cite where they
matter.
